\subsubsection{1H. Q4 / fourth-moment route: status and warning}\label{h.-q4-fourth-moment-route-status-and-warning}

A proposed minimal lemma Q4(\(d\)) asks whether, for degree-\(\le d\) \(p\) under Dobrushin,

\[\|p-\mathbb E_\mu p\|_{L^4(\mu)}
\le C_d\|p-\mathbb E_\mu p\|_{L^2(\mu)},\]
with \(C_d\) independent of \(n\) and \(p\).

\textbf{It is unproved, as are the derivative/ANOVA variants and the general induction.} The proposed Riesz estimate \(\|p-P_{t_*}p\|_4\le\frac12\|p\|_4\) was only a sketch and is not supplied. Conditional expectations can increase degree, so an induction cannot freely assume preservation.

If Q4 held at all degrees and with all derivative variants, then OZ's algebraic sup-norm bound plus Dobrushin Poincaré and reverse Hölder would give the displayed AA bound. But this does not constitute a simplification: Q4 encapsulates a very strong hypercontractive/DFKO-type statement and may be as difficult as the main target. Bonami and nonhomogeneous OZ constants carry exponential \(d\) factors, so calling the \(v^2\) result ``sharper'' without tracking them would be misleading.
